\chapter{SWHID syntax cheat-sheet}
\label{app:swhid-syntax}

This appendix summarises the \textsc{swhid} grammar. It is \emph{not} normative: the authority is
the \emph{SWHID Specification, version 1.2}~\autocite{swhidspec}, published openly on
swhid.org and standardised, word for word, as ISO/IEC~18670:2025~\autocite{isoiec18670}. The
section pointers below refer to that specification.

\section{Core identifiers (spec \S5)}

A \emph{core} \textsc{swhid} is four colon-separated fields,
\swhid{swh:1:\textit{type}:\textit{hash}}\,:\footnote{\url{https://www.swhid.org/specification/v1.2/5.Core_identifiers/}}
the scheme \texttt{swh}; the schema version (\texttt{1}); a three-letter object \emph{type}; and a
40-digit, lowercase-hex object \emph{identifier}. That identifier is a \textsc{sha}-1 computed over a
git-style serialisation of the object --- so a revision \textsc{swhid} is exactly the Git commit
hash with \texttt{swh:1:rev:} prepended. The five types are:

\begin{table}[htbp]
  \centering\small
  \begin{tabularx}{\linewidth}{@{}llX@{}}
    \toprule
    \textbf{Code} & \textbf{Object} & \textbf{Identifies} \\
    \midrule
    \texttt{cnt} & content   & the bytes of a single file \\
    \texttt{dir} & directory & a directory tree (names, permissions, and child hashes) \\
    \texttt{rev} & revision  & a commit, with its parents, author, and date \\
    \texttt{rel} & release   & a named release (annotated tag) \\
    \texttt{snp} & snapshot  & a repository's full state at one visit (all branches and tags) \\
    \bottomrule
  \end{tabularx}
\end{table}

\section{Qualified identifiers (spec \S6)}

A \emph{qualified} \textsc{swhid} appends a sequence of qualifiers to a core identifier, each
introduced by \texttt{;} and written \texttt{name=value}.\footnote{\url{https://www.swhid.org/specification/v1.2/6.Qualified_identifiers/}}
The core \emph{identifies}; the qualifiers \emph{locate} and \emph{contextualise}.

\begin{table}[htbp]
  \centering\small
  \begin{tabularx}{\linewidth}{@{}llX@{}}
    \toprule
    \textbf{Qualifier} & \textbf{Kind} & \textbf{Meaning} \\
    \midrule
    \texttt{origin} & context  & the URI the code was harvested from \\
    \texttt{visit}  & context  & the snapshot \textsc{swhid} of the visit \\
    \texttt{anchor} & context  & a node (\texttt{dir}/\texttt{rev}/\texttt{rel}/\texttt{snp}) the \texttt{path} is rooted at \\
    \texttt{path}   & context  & an absolute path from the anchor's root directory \\
    \texttt{lines}  & fragment & a line range within a \texttt{cnt} (1-based, inclusive; \texttt{A} or \texttt{A-B}) \\
    \texttt{bytes}  & fragment & a byte range within a \texttt{cnt} (0-based, inclusive) \\
    \bottomrule
  \end{tabularx}
\end{table}

The rules (spec \S6): each qualifier appears at most once; \texttt{visit} requires \texttt{origin};
\texttt{anchor} requires \texttt{path}; and the fragment qualifiers \texttt{lines}/\texttt{bytes}
apply only to a \texttt{cnt}. The recommended order is \texttt{origin}, \texttt{visit},
\texttt{anchor}, \texttt{path}, then \texttt{lines} or \texttt{bytes}.

\section{A worked example}

The specification's own example pins lines 9--15 of one file, in one revision, of one archived
repository:

\begin{verbatim}
swh:1:cnt:4d99d2d18326621ccdd70f5ea66c2e2ac236ad8b;
  origin=https://gitorious.org/ocamlp3l/ocamlp3l_cvs.git;
  visit=swh:1:snp:d7f1b9eb7ccb596c2622c4780febaa02549830f9;
  anchor=swh:1:rev:2db189928c94d62a3b4757b3eec68f0a4d4113f0;
  path=/Examples/SimpleFarm/simplefarm.ml;
  lines=9-15
\end{verbatim}

\noindent(written on a single line in practice). To compute or verify a \textsc{swhid}, see
Appendix~\ref{app:swhid-tools}.
